(1) $I$ commutes with every even element and $I^{2} = -1$ in $\Cl(1,3)$, so $(I\sigma_r)^{2} = I^{2}\sigma_r^{2} = -1$ using $\sigma_r^{2} = +1$. The eigenplane statement is \cite[Thm.~6.1]{ButlerT5} (\claimref{singularity.curvature-gauge-field-exact-shape-decoupling}), whose spectrum is $(-2,-2,1,1,1,1)$ with the doubled eigenvalue attained exactly on $\operatorname{span}\{\sigma_r, I\sigma_r\}$. (2) With $(I\sigma_r)^{2} = -1$ the exponential series splits into the cosine and sine series, giving the displayed rotor; $R_{2\pi} = \cos\pi + I\sigma_r\sin\pi = -1$, and conjugation by $-1$ is the identity, so the conjugation action closes at $2\pi$ while the rotor itself closes at $4\pi$. (3) The two flows are the two real one-parameter subgroups of the rotor group of the plane pair, related by the substitution $\lambda = i\theta$ in the exponential; the substitution exchanges $\cosh, \sinh$ with $\cos, \sin$ exactly because the two generators differ by the factor $I$ with $I^{2} = -1$. (4) Read with the conformal spinors' unit $I_{6} = I_{4}E$ as the unit of the continuation, the continued boost is $\exp(\theta I_{6}\sigma_r/2) = \tfrac12(1+E)R_\theta + \tfrac12(1-E)R_{-\theta}$, which is Proposition \claimref{synthesis.wick-rotor-chiral-angular-rotor}(ii): on the chiral half $E = +1$ it is $R_\theta$ and on $E = -1$ it is $R_{-\theta}$, so on spinors the Euclidean rotation is carried by the second member, with the sense fixed by the chirality. The restriction to spinors is forced: by Proposition \claimref{singularity.euclidean-rotor-where-it-lies}(3) the rotor $R_\theta$ fixes $e_0$, $e_1$ and rotates $e_2$ into $e_3$, along the bifurcation sphere, and no real rotor of $\Cl(1,3)$ realises the rotation around it, while the continued element is no rotor on vectors (Proposition \claimref{synthesis.wick-rotor-chiral-angular-rotor}(iii)).
